\begin{abstract}
Development tools report \emph{what} type an expression has, but not \emph{why} it has that type. This paper develops a theory of \emph{type slicing} that answers such questions: a programmer selects a term, queries any part of the type information associated with it, and receives a slice of the program---a well-formed partial program with irrelevant sub-terms folded away---that suffices to reproduce the queried type. We formulate type slicing for bidirectional type systems, where \emph{synthesis slices} explain the type a term synthesises and \emph{analysis slices} explain the type its surrounding context expects. The theory requires no cast dynamics as it applies to any bidirectional system equipped with a precision order on types and terms satisfying a downwards static graduality property. We develop the metatheory over a core calculus with holes, products, sums, and explicit polymorphism, based on the Hazelnut and marked lambda calculi, proving that every query has a minimal slice, and that refining a query monotonically shrinks its minimal slices. We then show how to calculate these slices both exactly and approximately. Finally, integrating type slicing with error marking theory extends these results to arbitrary ill-typed programs, so a single mechanism explains both types and type errors in \textit{complete}, \textit{incomplete}, and \textit{erroneous} code. The metatheory is mechanised in Agda, and a linear-time approximation of type slicing is implemented for the Hazel programming environment.
\end{abstract}

%%% Local Variables:
%%% mode: LaTeX
%%% TeX-master: "PolymorphicTypeSlicing"
%%% End:
